%! root=./main.tex
\section{Introduction}

% NOTE (@kjee): 1 page + 1/3 ~ 2/3 pages with about 3 contribution bullets.

* \cfi has been the cornerstone defense against various attacks launching control hijacking attacks targeting critical infrastructures.
  - While effectiveness in guarding control-sensitive memory regions against broader range of attacks launched by powerful and dedicated attackers, \cfi implements low-overhead (near-native) runtime overhead \wc{efficient} security defense.
  - Hence, \cfi has generated a significant body of research and has been widely adopted and deployed in production environments. \kjee{Even hardware-assisted solutions.}

* Despite its solid/promising/confirmed conceptual and theoretical foundations \wc{— which provide a reliable basis for protecting critical computing segments pervasive across various social and economic contexts in an efficient and effective manner —} the core security and safety questions regarding its actual implementations and real-world deployment are rarely asked and/or properly answered.

* A typical pain-point/loophole/challenges of general software development process still applies to the \cfi implementation
* As we have seen thus far beyond (not only for the) the case of \cfi and for the case of , manifesting a paper design (conceptual software design) into working, deployable software is not at all a trivial effort.
  - It generally requires significant development cost and many person-hours; Without any doubts, software development and implementation effort is an expensive and pain-staking task.
  - Even after implementation of a software prototype,  deployable to production, a long-lasting and tedious process of testing and patching remains necessary.
  - Several past incidents only have testified that software components with sound, verified designs can be implemented into prototypes bearing many bugs — some of them exploitable vulnerabilities with critical/deadly consequences — discovered only in later stages of testing.
  - Worse, some (few of them) vulnerabilities survive even comprehensive and rigorous process of testing, only to be discovered later with consequential impact.
    - Assurance via testing is fundamentally limited.

\kjee{we may want mention/state here that formal verification is an alternative to the solution; not necessarily going in deeper with complaints.}

* Such drawbacks/loopholes/limitations is considered to be more deadly/consequential/catastrophic for the case of security solutions as it critically undermines the trusted computing basis.
  - therefore we would like to say that the security assurance of the security defense solutions such as \cfi is utterly critical not only for its basic design but also for the software development and implementation processes.
  - however, we are yet to see proper solution to such problem.

* This fundamental gap between conceptual design and real-world software implementation \wc{artifact} is especially critical for security defense software: it gives users a false sense of assurance — making them less vigilant — while simultaneously introducing critical backdoors that allow an adversary to bypass defenses and compromise their target. \kjee{The SolarWinds case testified that exploitable vulnerabilities in security-sensitive infrastructure software are particularly deadly and consequential.}
  - The community's response to this critical blindspot has been largely limited to dynamic software testing. Despite significant efforts — \eg fuzzing, symbolic execution, and model checking — these approaches only improve coverage and reduce the likelihood of finding exploitable vulnerabilities; they are fundamentally limited in providing the level of assurance the community requires.

  \kjee{I want to come up with a proper phrasing of such phenomena.}

%%%%%%%%%%%%%%%%%%%%%%%%%%%%%%%%%%%%%%%%%%%%%%%%%%%%%%%%%%%%%%%%%%%%%%%%%%%%%%%%%%%%%%%%%%%%%%
% Now FM story begins; FM to the rescue.
%%%%%%%%%%%%%%%%%%%%%%%%%%%%%%%%%%%%%%%%%%%%%%%%%%%%%%%%%%%%%%%%%%%%%%%%%%%%%%%%%%%%%%%%%%%%%%
- A family of formal verification (machine check) approaches can be consider to be a solution/an answer.




{%
    % we want to motivate the reader so that they agree that formal verification of an implementation of binary rewriting is the best way forward
    % we need to convince them that people have tried the workarounds and that they just don't provide the same level of assurance that we will

    Development of mission-critical software systems, such as secure OSes, VMs, hypervisors, and runtime systems intended for deployment in adversarial environments, demands exceptionally high assurance development tools and procedures to meet strict safety, security, and correctness standards.
    Machine-checked formal verification has emerged as a hallmark approach for achieving the requisite guarantees for such systems,
    with success stories such as CompCert~\cite{compcert}, seL4~\cite{sel4}, and CakeML~\cite{cakeml} demonstrating the efficacy of such \emph{formal methods} (FM) for obtaining real-world, high-assurance systems in practice.
    These successes rely on mature theorem proving environments such as Rocq~\cite{rocq} and Isabelle/HOL, which have enabled rigorous, machine-checked reasoning about source-level semantics and compiler correctness.

    % \emph{Binary rewriters} transform binary programs to optimize, instrument, or enforce a security policy on them.
    % Rewriters are difficult to verify due to their complexity and the lack of tooling capable of formally reasoning about underspecified pieces of code.
    % The accepted solution to make a rewriter trustworthy is to split the task of creating a sound rewritten binary across two pieces of software.
    % The untrusted rewriter rewrites the target program without guaranteeing its output is a sound transformation.
    % The formally verified \emph{verifier} checks the output and accepts it if it is correct.

    While machine-certified compilation has received extensive attention in the literature (cf., \cite{lustre_compiler,certi_crypto_comp,certicoq_wasm,compcerto,rocq_compilers_alacarte,design_impl_certifying_compiler,certified_choreographies}), binaries that execute on real mission-critical systems are seldom the unaltered output of a compiler.
    Compiled binaries typically undergo a variety of binary-level transformations prior to execution, including load-time linking, hardening, and sandboxing, which prepares the binary to safely execute in the context of the surrounding OS or VM\@.
    The \emph{binary rewriters} that implement these transformations demand an equivalent or higher level of assurance, since any flaw could undermine all assurances achieved during compilation.
    However, complete formal verification of a binary rewriter \emph{implementation} has remained a difficult challenge for FM, primarily because their correctness properties are universal quantifications over arbitrary binary code.
    In particular, a binary rewriter is correct if and only if it achieves a desired safety property for all possible input binaries, not merely those produced by a particular compiler or that limit themselves to a particular subset of the ISA\@.

    Binary \emph{control-flow integrity} (CFI) enforcement is one particularly prominent example of the associated challenges~\cite{safedispatch,typearmor,object_flow_integrity,fineibt,renewablejitcfi,eccut,resource_constrained_cfi,fixingfineibt}.
    Source-agnostic CFI algorithms transform input binaries so as to confine all their reachable control-flows to a policy-specified whitelist of edges.
    Implementations typically identify and instrument the target binary's control-flow transfer instructions (e.g., jumps, calls, and returns) with extra \emph{guard code} that checks each impending destination against the policy at runtime.
    Impending violations solicit a safe intervention, such as secure abort, rollback, or restart.
    A binary CFI implementation is \emph{sound} if all reachable execution traces of all binaries it generates obey the policy.

    Since source-driven (\emph{top-down}) FM approaches cannot state or verify a source-agnostic soundness property, and since binary-driven (\emph{bottom-up}) FM approaches are typically designed to verify concrete input binaries (not the arbitrary output of a binary transformation), prior efforts to verify binary CFI have resorted to verifying each CFI-hardened binary case-by-case (e.g., \cite{rocksalt}).
    This helps code-consumers protect themselves from untrusted code-producers, but it does not empower code-producers to create a verifiably sound CFI implementation that always passes the consumer's checks.
    Producers must resort to unit testing (i.e., running the rewriter on many input binaries to see whether any resulting binaries fail verification), losing the primary benefit of FM\@.

    Formally verifying a binary rewriter implementation requires overcoming two major hurdles:
    First, it requires a verification framework that supports \emph{underconstrained-code properties}---properties in which the target code is a quantified proof variable, not a concrete program.
    Such properties are inexpressible in most existing bottom-up FM frameworks (e.g., \cite{trabin,isabil,islaris,myreen08,bjorn14}).
    Second, the approach must have a capability for reasoning about arbitrary byte sequences as (possibly invalid) binary code.
    This is necessary because rewriter \emph{implementations} read and write bytes, and their soundness is therefore contingent upon their correct interpretation of those bytes as instructions.
    Idealizing a rewriter as an abstract model of its algorithm waives verification of many of the most complex parts of its implementation, which are pernicious sources of potential error.

    %%% KWH Comment: The relationship of this work to ours is too complex for the intro, imo.  We should discuss this in more depth in the related work section.
    %The only work in the literature which may support this was released in 2025 and focused on relational properties of \emph{concrete} programs with statically known code~\cite{relational_realistic}.


    % Verifying binary rewriters requires reasoning about their effects on arbitrary input codes.
    % Historically, binary analysis frameworks have been limited to analyzing concrete code~\cite{bap,isabil,islaris,myreen08,bjorn14,trabin}, preventing them from reasoning about binary transformations.
    % In 2025 a new binary analysis platform debuted with this capability but did not demonstrate its application, focusing instead on relational logic~\cite{relational_realistic}.

    % \emph{Top-down FM} yields high assurance for software representable as source code by extending proofs of source-level properties down to executable codes, such as through type-preserving certified compilation.
    % But codes that are binary-only, and proofs of binary-level properties that are independent of sources, require a different approach.
    % For example, low-level runtime ABIs often contain raw assembly or machine code that has been carefully optimized for maximal performance on specific instruction set architectures (ISAs).
    % These optimizations benefit all source codes that link to them, but potentially undermine top-down FM guarantees that only prove properties of the compiler-generated parts of the final executable, not the ABI code contributed by the runtime system during linking.
    % Likewise, the correctness of software that manipulates binary code, such as VMs, process loaders, linkers, and binary security hardening transforms, cannot be expressed purely as source-level properties because they concern software that takes \emph{arbitrary} binary code as an \emph{input}.
    % For example, a VM implementation is secure only if it can successfully load and execute arbitrary byte sequences (even those not generated by any compiler) in a way that prevents the loaded program from escaping the VM sandbox.

    % \emph{Bottom-up FM} verifies binary code by lifting~\cite{formally_verified_lifting} it to a low-level intermediate language (IL) that expresses the semantics of machine instructions as cpu state transitions.
    % When the IL has a formally defined semantics, it is amenable to interactive theorem proving, facilitating the same style of compositional, mechanized reasoning used in top-down FM, but applied directly to executable code.
    % This capability affords stating and proving properties of all possible executions.

    % But while existing bottom-up FM frameworks facilitate reasoning about concrete binary codes, the task of formally reasoning about arbitrary binary codes and the programs that read them has remained an unsolved challenge.
    % As a hallmark example, \emph{binary control-flow integrity} (CFI) algorithms~\textcolor{red}{[citations]} instrument input binaries to constrain their control-flows to a whitelist of permissible control-flow transfer edges.
    % The output of the CFI transform is a self-monitoring, self-securing binary that requires no additional sandboxing.
    % CFI security is a core technology for enforcing many critical software security policies, including process isolation, software fault isolation, stack integrity, mobile code security, and code-reuse attack detection and defense~\textcolor{red}{[citations throughout this list]}.

    % Although prior efforts have verified idealized binary CFI \emph{algorithms} for abstract ISA models of varying detail~\textcolor{red}{[citations]}, % no formally verified binary CFI \emph{implementation} for a real, mainstream ISA has yet been realized.
    % none have yet verified a binary CFI \emph{implementation} for a real mainstream ISA.
    % This has left open the dangerous possibility that binary CFI implementations might not match the algorithms on which their designs are informally based.
    % In practice, such implementations have many subtle details that are prone to error, including accurate decoding of all reachable instructions in potentially malicious or malformed binaries, correct identification of all possible indirect control-flow transfer operations, complete mediation of all such operations, and careful rebasing and relinking of all code fragments displaced by guard code introduced during the rewriting process.
    % A full proof of safety for a binary CFI implementation must therefore consider all (arbitrary) input byte sequences that constitute executable code on the ISA, and prove that all possible executions of the implementation's output byte sequence satisfy the control-flow policy.

    %Although checkers provide assurance when they accept a binary, they offer no guarantees or guidance when they reject one, leaving the user uncertain about the safety of the transformed code.

    In this paper we address this long-standing challenge through the design and implementation of \Picinae{}, a \uline{p}latform \uline{i}n Ro\uline{c}q for \uline{in}struction-level \uline{a}nalysis of \uline{e}xecutables.
    \Picinae{} is designed from the ground up to construct machine-checked proofs about real binary programs for real ISAs, while maintaining sufficiently general foundations to prove theorems in which all or part of the target binary code is a proof variable rather than a specific program.
    It supports modeling both von Neumann and Harvard architectures, including difficult ISA features like misaligned instructions on Intel, conditional instruction prefixes on ARM, undefined and non-deterministic effects, memory access permissions, self-modifying code, hardware exceptions, indirect calls and returns, mixed endianness on ARM, instructions with internal microcode loops, and multi-ISA integrated systems.
    These capabilities are leveraged to prove that a realistic binary CFI implementation for a mainstream ISA (ARM7) safely hardens arbitrary binaries (including rejecting unsafe or malformed binaries that it cannot secure).
    Specifically, we present the following contributions:
}

\begin{itemize} [topsep=0pt,itemsep=-1ex,partopsep=1ex,parsep=1ex, leftmargin=*]
    \item The \emph{first} binary-agnostic formal verification framework the most significant/foundational binary inline monitor security solution, \cfi.
    % \item the \Picinae{} framework, which facilitates machine-checked proofs of correctness for binary programs and binary transformations;
    % \item a dependently-typed, co-inductive foundation for proving correctness properties of algorithms that manipulate machine code; and
    % \item a case study demonstrating the first machine-checked proof of soundness for a binary CFI rewriter implementation targeting ARM7 microprocessor code.
\end{itemize}
